Der Hang ist ein rechtwinkliges Dreieck: $\ell = \lP$ entlang des Hangs, $h = \hP$ nach unten, also waagrecht (Pythagoras)
\begin{align}
 b &= \bPF = \bP
\end{align}
Für den Neigungswinkel $\alpha$ gilt $\sin\alpha = h/\ell$ und $\cos\alpha = b/\ell$. Die Gewichtskraft wird zerlegt: $m\,g\sin\alpha$ hangabwärts, $m\,g\cos\alpha$ in den Hang hinein.

\begin{abcliste}
\abc
Die Normalkraft hält der Komponente in den Hang hinein das Gleichgewicht:
\begin{align}
 \ssc{F}{N} &= \FNF \\
   &= \m\cdot\g\cdot\frac{\bP}{\lP} \\
   &= \FN \approx \result{\FNP{0}{2}}
\end{align}

\abc
Entlang des Hangs wirken die Hangabtriebskraft und die Reibung $\mu\,\ssc{F}{N}$:
\begin{align}
 m\,a &= m\,g\sin\alpha - \mu\,m\,g\cos\alpha \\
 a &= \aF \\
   &= \g\cdot\left(\frac{\hP}{\lP}-\muG\cdot\frac{\bP}{\lP}\right) \\
   &= \a \approx \result{\aP{0}{2}}
\end{align}
Die Masse kürzt sich weg: Ein schwerer und ein leichter Skifahrer werden gleich schnell schneller (ohne Luftwiderstand).
\end{abcliste}
